\documentclass[11pt]{article}
\usepackage[letterpaper,margin=1in]{geometry}
\usepackage[T1]{fontenc}
\usepackage{lmodern,amsmath,amssymb,amsthm,booktabs,longtable,array,hyperref,xcolor}
\hypersetup{colorlinks=true,linkcolor=blue,urlcolor=blue}
\newtheorem{theorem}{Theorem}[section]
\newtheorem{lemma}[theorem]{Lemma}
\newtheorem{proposition}[theorem]{Proposition}
\newtheorem{corollary}[theorem]{Corollary}
\theoremstyle{definition}\newtheorem{definition}[theorem]{Definition}
\theoremstyle{remark}\newtheorem{remark}[theorem]{Remark}
\newcommand{\Z}{\mathbb Z}\newcommand{\N}{\mathbb N}
\newcommand{\Con}{\operatorname{Con}}
\newcommand{\height}{\operatorname{ht}}
\setlength{\emergencystretch}{3em}
\title{Residual basis reduction and smaller Diophantine consistency equations}
\author{Research prepared for Vladimir Reshetnikov}
\date{25 October 2026}
\begin{document}
\maketitle
\begin{abstract}
This report studies explicit polynomial equations whose unsolvability is
provably equivalent in ZFC to the syntactic consistency of ZFC. The main
new method changes the basis of quadratic residual equations to reveal
linear consequences, then eliminates their unit-coefficient coordinates.
It reduces both the dimension and the coefficient cost of a direct
proof-calculus construction. Further results combine a shift-free
descending-factorial certificate with residual reduction, replace
unnecessarily large Pell transport denominators, and remove redundant
arithmetic gates from small-variable root certificates.
All comparisons specify the domain and the measured cost. The unbounded
arguments use explicitly retained ordinary arithmetic and refutation
calibrations. Exact finite checks certify the supplied polynomials and
transformations; they are not a new end-to-end proof-assistant certificate,
a consistency proof, or a theorem of global optimality.
\end{abstract}
\tableofcontents
\section{The problem and the comparison boundary}
For a fixed integer polynomial $P$ and a stated domain $D$, the target is
\[
 \mathrm{ZFC}\vdash
 \bigl(\neg\exists \boldsymbol{x}\in D^n\;P(\boldsymbol{x})=0\bigr)
 \longleftrightarrow\Con(\mathrm{ZFC}).
\]
Here consistency means absence of a finite syntactic refutation in the
specified effective proof system. The statement is obtained by proving
solvability equivalent to existence of such a refutation and then negating
both sides. It does not decide whether either side holds.

There is no single total order of simplicity. We measure the number of
quantified variables, their domain, total degree, number of quadratic
residuals where applicable, support of the fully collected polynomial,
maximum absolute integer coefficient, and live binary operations in a
specified straight-line evaluator. Addition, subtraction and
multiplication each cost one, including squares and multiplication by
fixed constants. Reusing a register is free. Fixed numerical ports are
not variables, but their entire coherent tuples must be specified.

The immutable repository comparison is
\texttt{e0c98c629bf10d30b695884f6b1d2ab361bdf513} in
\href{https://github.com/VladimirReshetnikov/ProveIt}{ProveIt}.
The archive retains the source identities and distinguishes actively
accepted repository constructions from newly submitted incoming packages.
In particular, the incoming \emph{Cofinal Anchors} and \emph{Certificate
Reductions} packages are stronger comparators than our preceding report.
Our reductions of their literal equations are explicit; the word
``incoming'' does not imply an independent external review or repository
adoption. Attribution and inherited mathematical hypotheses are preserved.

\subsection{Selected results}
\begin{center}\small
\begin{tabular}{lrrrrr}\toprule
Domain and variant & Variables & Degree & Rows & Terms & Height\\\midrule
$\Z$, dimension &278&4&188&12884&12227\\
$\Z$, balanced &279&4&189&5483&7200\\
$\Z$, sparse &279&4&189&5202&12662\\
$\Z$, fewer terms &284&4&194&4270&10864\\
$\Z$, sparsest listed &306&4&216&3870&14512\\\midrule
$\N$, balanced &229&4&179&5607&65892965\\
$\N$, smaller lower height &228&4&178&9522&560804902\\
$\N$, smaller sparse &228&4&178&7512&2445412949744\\\bottomrule
\end{tabular}
\end{center}
Here $\N=\{0,1,2,\ldots\}$ and $\Z$ denotes unrestricted integers.
The 279-variable balanced equation improves the incoming 282-variable
comparison on every reported varying measure: its parent has 192 rows,
5491 terms and height 128446. The 278-variable equation goes one variable
further at a support cost. The other integer comparisons are to the
matched 287-variable/197-row/4313-term and
309-variable/219-row/3922-term incoming equations, both of height 128446.
These are direct proof-calculus reductions, not absolute variable-count
records across all coefficient encodings.

\begin{center}
\begin{tabular}{lrrr}\toprule
Domain and branch & Variables & Exact degree & Binary operations\\\midrule
Strictly positive, $D=XY$, squared weight &10&626&77\\
Strictly positive, $D=XY$, product weight &10&464&78\\
Strictly positive, $D=q^4$, product weight &10&460&79\\
Strictly positive, partial normalization &10&384&80\\\midrule
$\Z$, auxiliary root and odd guard &14&86&112\\
$\Z$, strong root and odd guard &14&94&110\\
$\Z$, exact recompiled predecessor &13&156&125\\\bottomrule
\end{tabular}
\end{center}
The new 14-variable degree-86 equation improves both degree and evaluator
cost over the preceding degree-88/119-operation form. Degree 94 is a
lower-cost companion, not a degree improvement. A degree-88/115-operation
exact recompilation is also retained as a check. The 13-variable entry
preserves the entire earlier polynomial and only shortens its evaluator.
No smaller fixed-port payload or fully expanded coefficient height is
claimed for these circuit branches.

\section{Changing residuals without changing solutions}
\begin{lemma}[Elementary residual operations]
Let $R=(R_1,\ldots,R_m)$ be integer polynomials. Replacing $R_i$ by
$R_i+kR_j$, for $i\ne j$ and $k\in\Z$, preserves their common zero set
on every fixed coordinate domain. Its inverse replaces the new $R_i$
by $R_i-kR_j$. Consequently, over $\Z^n$ or $\N^n$, the sums of the
squares of the old and new residuals have precisely the same zeros.
\end{lemma}
\begin{proof}
The forward and inverse substitutions are polynomial identities and
leave all other rows unchanged. A sum of integer squares vanishes if
and only if every summand vanishes. The assertion about sums of squares
uses this ordered-domain fact; it is not asserted over the complex numbers.
\end{proof}

\begin{lemma}[Exposing an affine coordinate]
Suppose exact elementary row operations produce a row
$\varepsilon x_j+L(\boldsymbol{x}_{\ne j})$, where
$\varepsilon\in\{1,-1\}$ and $L$ is affine with integer coefficients.
Substitute $x_j=-\varepsilon L$ in every row and delete the resulting
zero row and the coordinate $x_j$. The remaining system has a bijective
correspondence of integer zeros with the original system, with the displayed
formula as inverse coordinate restoration. Quadratic residuals remain
quadratic. The assertion for natural-number coordinates additionally
requires a proof that the restored coordinate is nonnegative.
\end{lemma}
\begin{proof}
At every zero the displayed row uniquely specifies $x_j$. Conversely,
substitution followed by restoration makes that row zero and makes
every retained row equal to its substituted counterpart. Affine
substitution does not increase degree. Repeat in the recorded order
for several eliminations and invert in reverse order.
\end{proof}

These elementary facts motivate a different optimization problem from
changing the encoded proof calculus or its Pell witnesses. Take the matrix
of quadratic coefficients of the residual list. Its integer row
combinations can have a vanishing quadratic part even when none of the
individual rows is affine. Such combinations expose otherwise hidden
coordinates. Afterwards a second residual-basis search can lower the
support or coefficient height of the sum of squares. The final polynomial
usually changes, while its zero set does not. Every accepted operation is
recorded and replayed exactly; heuristic choices made during the search
are not part of the proof.

\section{Hidden affine consequences: quartics in 278 and 279 integers}
\label{sec:row-exposure}

The new direct-calculus contribution is a reduction of the incoming
282-coordinate quartic to 279 coordinates without a support or height
penalty, and to 278 coordinates at a stated support tradeoff.  The mechanism
is not a new number-theoretic assumption: elementary changes of residual
basis expose affine identities that the original gate presentation hides.
The corresponding 287- and 309-coordinate parents also lose three
coordinates.  Every coordinate remains an unrestricted integer.

\begin{center}\small
\begin{tabular}{lrrrr}
\hline
Equation & Variables & Rows & Monomials & Height\\
\hline
Incoming dimension parent &282&192&5491&128446\\
Incoming grouped parent &287&197&4313&128446\\
Incoming sparse parent &309&219&3922&128446\\
\hline
New dimension tradeoff &278&188&12884&12227\\
New balanced &279&189&5483&7200\\
New sparse 279 &279&189&5202&12662\\
New sparse 284 &284&194&4270&10864\\
New sparse 306 &306&216&3870&14512\\
\hline
\end{tabular}
\end{center}

All eight displayed equations have exact degree four.  The two
279-coordinate equations improve every varying metric of the matched
282-coordinate parent; the 284- and 306-coordinate equations do likewise
for their matched parents.  The 278-coordinate point sacrifices support
for one fewer variable.  These are literal finite endpoints, not
global-minimality or literature-priority claims.

\subsection{Equivalence lemmas}

\paragraph{Residual basis.}
If $i\ne j$ and $k\in\mathbb Z$, replacing $R_i$ by $R_i+kR_j$
leaves the common zero set unchanged: one implication is substitution,
and the other follows by subtracting $kR_j$ after using $R_j=0$.
The inverse uses $-k$, so each such operation is unimodular.  The
sum-of-squares polynomial generally changes; its zero set does not.

\paragraph{Integral coordinate chart.}
The substitution $x_t=z_t+L(z)$, with $L$ integral affine and independent
of $z_t$, is a bijection of integer lattices.  Its inverse is
$z_t=x_t-L(x)$.  A sequence is inverted in reverse order.  Quadratic
degree is preserved.

\paragraph{Unit affine projection.}
If one residual is $\epsilon x_t+L=0$, with $\epsilon\in\{1,-1\}$
and $L$ integral affine and independent of $x_t$, substitute
$x_t=-\epsilon L$ into every other residual and delete this row and
coordinate.  Restriction maps old zeros to new zeros; the displayed
integral formula gives the unique inverse extension.  Repeated extensions
are evaluated in reverse projection order.  Degree remains at most two.

Since $\sum_i R_i^2=0$ over $\mathbb Z$ exactly when every $R_i=0$,
these three facts prove the corresponding equivalence for a single quartic.
No sign guard, divisibility condition or cofinal choice is introduced.

\subsection{Credited preconditioning and the new identities}

Six previously published-in-repository coordinate shears are transferred
to the three incoming parents.  This preconditioning is credited prior
work, not the novelty claim.  With zero-based coordinate indices, the
successive old-coordinate substitutions are
\[
\begin{aligned}
x_{21}&\leftarrow x_{21}-2x_{10}-2x_{11},&
x_{23}&\leftarrow x_{23}-8x_{22},\\
x_{38}&\leftarrow x_{38}-8x_{37},&
x_{35}&\leftarrow x_{35}-4x_{10}-5x_{12}-2x_{25}-x_{27},\\
x_{27}&\leftarrow x_{27}-x_{10}-x_{12},&
x_{25}&\leftarrow x_{25}-x_{11}-x_{12}.
\end{aligned}
\]
They are executed in the displayed order, not simultaneously.  The
upstream certificates, exact source hashes and commit are retained.
Each forward and inverse map is independently checked here; no semantic
acceptance claim for the upstream V-R1 endpoint is needed.

Let $Q_i$ denote the residuals after this preconditioning.  Exact
coefficient collection in each of the three parents gives
\[
\begin{aligned}
Q_{21}-Q_6
 &=4+4x_{10}+8x_{11}+4x_{21}-x_{22}+x_{37},\\
Q_{73}-Q_{71}-4Q_{70}-8Q_{68}
 &=12-6x_3-8x_{84}-4x_{85}-x_{86}+x_{89},\\
Q_{122}-Q_{120}
 &=3-2x_{107}+2x_{151}-x_{167}+x_{169}.
\end{aligned}
\]
These are polynomial identities, not consequences inferred from sample
solutions.  Replace rows 21, 73 and 122 by the left sides, using five
elementary row shears.  The donor rows are unchanged.  The unit-affine
projection lemma then gives
\[
\begin{aligned}
x_{37}&=x_{22}-4-4x_{10}-8x_{11}-4x_{21},\\
x_{89}&=x_{86}-12+6x_3+8x_{84}+4x_{85},\\
x_{169}&=x_{167}-3+2x_{107}-2x_{151}.
\end{aligned}
\]
The deleted coordinates are named \texttt{add.145}, \texttt{mul.346}
and \texttt{base4.T}.  This proves the drops
$(282,192)\to(279,189)$, $(287,197)\to(284,194)$ and
$(309,219)\to(306,216)$ in coordinates and rows.

\subsection{One further projection and the 278-coordinate point}

The already affine row $Q_{90}$ has unit coordinate $x_{105}$,
named \texttt{add.380}.  It is not counted as a newly discovered affine
identity.  Its unique restoration is
\[
\begin{aligned}
x_{105}={}&x_{10}+x_{11}+4x_{21}+2x_{24}+4x_{26}+2x_{29}
            +2x_{31}+4x_{32}\\
          &+2x_{34}+2x_{36}+4x_{39}+2x_{40}+2x_{43}
            +2x_{47}+4x_{68}+2x_{69}.
\end{aligned}
\]
Applying this additional projection to the dimension parent gives 278
coordinates and 188 rows.  The dense replacement explains the larger
collected support.  No new semantic lemma is needed.

\subsection{The fixed quartics and their logical meaning}

Each endpoint is defined exactly by
\[
 P(\mathbf z)=\sum_{i=0}^{m-1}r_i(\mathbf z)^2,
 \qquad \mathbf z\in\mathbb Z^n,
\]
where its complete list of integral quadratic $r_i$ is in
\texttt{equation.json}; complete collected coefficients are in
\texttt{expanded\_polynomial.json}.  The package also supplies every row
as human-readable text and TeX, the full coordinate-order table, and a
complete affine map restoring the original parent coordinates.

Following the projections, a frozen list of elementary residual row
shears selects each stated support/height profile.  The lists are fully
explicit in the endpoint certificates.  Their proof is the residual-basis
lemma applied successively, not the search heuristic that found them.
The deterministic builder runs no optimizer.  Exact expansion confirms
degree four and that every declared coordinate occurs.

\paragraph{Theorem.}
Each supplied new equation has an integer zero if and only if its
matched incoming parent has an integer zero.  Hence, composing with the
parent's ordinary C222/ZF/ZFC refutation calibration, ZFC proves that
integer unsolvability of each new equation is equivalent to
$\operatorname{Con}(\mathrm{ZFC})$.

\paragraph{Proof.}
For a new zero, undo the final row shears to recover all projected
residual equations.  Restore the deleted coordinates in reverse order;
the deleted equations vanish identically and all remaining equations
hold.  Undo the affine-exposure row shears, then evaluate the coordinate
charts in reverse execution order.  This produces an old zero.  In the
other direction, use the inverse coordinate bijection, the exact affine
equalities, restriction to retained coordinates, and the forward row
shears.  Every step is integral and reversible on zeros.  The final
logical equivalence is exactly the inherited parent calibration.
This argument does not prove consistency or claim a new end-to-end
kernel-checked Lean theorem.

\subsection{Verification and the limit of this particular method}

An independent standard-library checker imports no producer.  It checks
coordinate inverses, elementary row inverses, unit coefficients, the
identical zero of every deleted row, absence of deleted coordinates,
the exact coordinate renumbering, every collected coefficient and metric,
and coordinate liveness.  It also reconstructs every original coordinate
as an affine polynomial in the final ones and checks every original
residual pullback as a polynomial identity.  Five corrupted-artifact
controls are rejected.  Normal and \texttt{python -O} runs pass, including
fresh byte-identical rebuilds.  The printed residual text and TeX are
parsed back to their complete coefficient lists.

For the 278-coordinate endpoint, the quadratic parts of the 188 residuals
have rank 188 modulo 1000003.  Thus they also have rank 188 over
$\mathbb Q$: a nonzero modular minor cannot have zero integer
determinant.  No nonzero rational linear combination of this final
residual basis is affine.  The parent's corresponding rank is 188, and
the four displayed independent affine combinations account for its
four-dimensional kernel.  This finite obstruction concerns this row
span only, not nonlinear ideal consequences or global variable
minimality.  The optional rank replay is supplied separately.

\section{Shift-free descending factorials and cofinal natural quartics}
\label{sec:shift-free-natural}

The new natural-domain branch combines the descending factorial identity
from the incoming \emph{Certificate Reductions} packet with the cofinal
semantic anchor and strict amplification exponent from Round~17. Its new
step removes the artificial positive shifts on both the factorial and the
CRT-carrier coordinates, deriving their positivity in the correct order
from the already present polynomial rows.

\begin{center}
\begin{tabular}{lrrrr}
\hline
Endpoint & Variables & Rows & Terms & Height\\
\hline
Round~17 balanced &229&179&9184&83954634\\
New projected precursor &229&179&6596&83879378\\
New balanced, row-reduced &229&179&5607&65892965\\
Round~17 smaller &228&178&9866&2445412949744\\
New smaller, lower height &228&178&9522&560804902\\
New smaller, sparser &228&178&7512&2445412949744\\
\hline
\end{tabular}
\end{center}
Every row concerns a quartic over independently nonnegative integers.
The two 228-variable variants use the identical new semantic graph and
only change the certified elimination order; all restorations are proved
nonnegative. In particular, no integer-only coordinate projection is
being applied to natural variables.

\subsection{The exact polynomial and logical statement}
The balanced polynomial is
\[
 P_{229}(x_0,\ldots,x_{228})=\sum_{j=0}^{178}R_j(x_0,\ldots,x_{228})^2,
 \qquad (x_0,\ldots,x_{228})\in\N^{229}.
\]
The full literal residual list is supplied in
\texttt{balanced229/residuals.tex} and \texttt{equation.json}; every collected
coefficient is in its \texttt{expanded\_polynomial.json}. The analogous files
in \texttt{low\_height228/} and \texttt{sparse228/} define their complete
178-row polynomials. A JSON term \([c,[i,j,\ldots]]\) means
\(c x_i x_j\cdots\); an empty index list denotes a constant. These
polynomials contain no binomial, factorial, floor, power, or Pell primitive.

With the retained ordinary numerical C222/ZF/ZFC refutation calibration,
each polynomial has a natural zero exactly when a finite syntactic ZFC
refutation exists. Consequently ZFC proves the equivalence between its
no-natural-zero sentence and \(\operatorname{Con}(\mathrm{ZFC})\).
This neither proves consistency nor supplies a new end-to-end kernel
proof or a global minimality theorem. The coefficient verification below
is finite; the unbounded arithmetic steps are proved separately.

\subsection{Semantic positivity replaces two coordinate shifts}
Keep \(t=t_0+12\), but now use unshifted natural coordinates \(F,n,B,r,s\)
and the three polynomial rows
\begin{equation}\label{eq:sf-certificate}
 n+1=F^2,\qquad V=BF+r,\qquad B=r+s+1.
\end{equation}
The first exact power request is \(V=n^t\). Its retained natural modulus
coordinate \(T_n\) obeys
\[
 T_n+n^2+1=2An.
\]
This literal equality excludes \(n=0\), even before any power or factorial
is interpreted. Thus \(F\ge2\) and \(n=F^2-1\ge3\).
The new row \(n+1=F^2\) is deliberately not marked as a defining gate
with globally nonnegative restoration \(F^2-1\).

The remaining exact requests, with the same self-sized radix as Round~17,
are
\begin{equation}\label{eq:sf-powers}
 V=n^t,\quad N=B^x,\quad q=(NS)^x,\quad
 D_t=q^t,\quad D_y=q^{y+1},\quad D_x=q^x,
\end{equation}
where
\[
 S=1+C+A_a+A_b+A_c+A_d+A_S+A_{S'}+A_u+A_v.
\]
All representatives are positive. The shared frame is the sum of every
distinct power base, exponent and output, together with every modular top
index, plus 15. There are thirteen nonanchor seeds and eleven modular
requests. The scale equations give \(A\ge(W+1)^W\), and the one-row
anchor described below authenticates all indices before the powers are
recovered. No factorial identity is needed for this stage: all atoms are
natural and all intended exponents are positive, since \(t\ge12\),
\(x\ge F+1\ge3\), and \(y\ge F+1\ge3\).

\begin{lemma}[Shift-free bootstrap]\label{lem:sf-bootstrap}
Once the exact powers are recovered, the rows
\eqref{eq:sf-certificate} imply
\[
 B>F\ge2,\qquad V<B(F+1)\le B^2\le N\le q.
\]
In particular, the first binomial can be decoded without assuming
\(F=t!\) or assuming that its output is positive in advance.
\end{lemma}
\begin{proof}
The remainder rows give \(B\ge1\) and \(V<B(F+1)\).
For \(F\ge2\),
\[
 n=F^2-1\ge F+1,\qquad V=n^t\ge n^2>F(F+1).
\]
Therefore \(B>F\), and integrality gives \(B\ge F+1\).
Since \(x\ge3\), \(B^2\le B^x=N\le(NS)^x=q\).
For \(0\le j\le t\),
\(\binom nj\le n^j\le n^t=V<q\), with \(\binom nj=0\) when
\(j>n\). Thus every relevant coefficient is carry-free.
\end{proof}
The root-modulus divisibility needed at this early point is also already
available. The raw endpoints are
\[
 y\ge F+1,\qquad x\ge F+1,\qquad
 x=1+(1+(t+1)y)u_t.
\]
As \(x\ge3\), \(u_t\ge1\), giving
\(x\ge2+(t+1)y>\max(t+1,y+1)\). Hence
\(M=Nq^x\) is divisible by \(qD_t\) and by all required range
denominators multiplied by \(N\), without yet knowing \(F=t!\).
The exact modular digit now gives \(B=\binom nt\). Since \(B\ge1\),
\(n\ge t\), and the two remainder rows say
\(F=\lfloor n^t/\binom nt\rfloor\).

\begin{lemma}[Descending fixed point]\label{lem:sf-fixed}
If \(t\ge3\), \(F\ge2\), and \(n=F^2-1\ge t\), then
\[
 F=\left\lfloor\frac{n^t}{\binom nt}\right\rfloor
 \quad\Longleftrightarrow\quad F=t!.
\]
\end{lemma}
\begin{proof}
Put \(T=t!\), \(U=t(t-1)/2\), and
\[
 H=\frac{n^t}{\binom nt}
   =T\prod_{j=0}^{t-1}(1-j/n)^{-1}\ge T.
\]
If \(F=\lfloor H\rfloor\), then \(F\ge T\).
The inequality \(T-1>U\) holds at \(t=3\) and persists by induction:
from \(T\ge U+2\) one gets
\((t+1)T-1\ge(t+1)(U+2)-1>U+t\).
Consequently
\(n\ge T^2-1=(T+1)(T-1)>(T+1)U\).
The elementary inequality \(\prod(1-a_j)\ge1-\sum a_j\) gives
\[
 T\le H\le\frac{Tn}{n-U}
   =T+\frac{TU}{n-U}<T+1.
\]
Thus \(F=T\). Conversely the same estimates at \(F=T\) prove that
\(\lfloor H\rfloor=T\).
\end{proof}

\subsection{The descending carrier and its correctly oriented index}
\begin{lemma}\label{lem:sf-crt}
Let \(F=t!\), \(t\ge3\). Then
\[
 d_i=F^2/i-1\quad(1\le i\le t)
\]
are pairwise coprime positive integers greater than one, every prime divisor
of each \(d_i\) exceeds \(t\), and
\(\prod_i d_i=\binom{F^2-1}{t}=B\).
\end{lemma}
\begin{proof}
Since \(i\mid F\), \(d_i\equiv-1\pmod F\), excluding all primes at
most \(t\). A common prime of \(d_i,d_j\) for \(i\ne j\) would divide
\(id_i-jd_j=j-i\), impossible for a prime exceeding \(t\).
The product identity is
\(\prod_i(F^2-i)/i=\binom{F^2-1}{t}\).
\end{proof}
The literal index equality is now
\begin{equation}\label{eq:sf-index}
 C+1=F^2+Bq_C,\qquad\text{equivalently } C=n+Bq_C,
\end{equation}
not the rising-carrier predecessor's differently signed row.
Modulo \(d_i\), it identifies \(C=i-1\).
Its quotient is natural in completeness: \(F\ge6\),
\(3\le t\le n-2\), and
\(B=\binom nt\ge\binom n2>n\). Every positive representative
congruent to \(n\pmod B\) is at least \(n\).

\subsection{The shared Pell anchor and recovery chronology}
For \(A\ge2\), define the ordinary Pell sequences by
\[
 (A+\sqrt{A^2-1})^j=X_j(A)+Y_j(A)\sqrt{A^2-1}.
\]
The complete retained three-norm certificate indexes a single principal
pair at the derived label \(J\). The simultaneous row is
\begin{equation}\label{eq:sf-anchor}
 Y_J(A)=J+\sum_i(Y_i-k_i)^2+(A-1)Q,\qquad Q\in\N.
\end{equation}
Here the notation \(Y_J(A)\) is its proved meaning, not a primitive of
the polynomial. Before indexing, the row supplies the anchor's omitted
inequality \(J\le Y_*\). Its full private witnesses remain natural.

For clarity, the retained cap argument applies whenever
\(W\ge3\), \(2\le J<W\), \(A-1>W^3\), at most \(W-2\)
seeds occur, and \(1\le k_i<W\).
Pell classification writes \(Y_i=Y_{j_i}(A)\). If \(j_i\ge J\),
then \(Y_i\ge Y_J(A)=V_*\ge2A>2W\), so
\((Y_i-k_i)^2\ge(V_*-W+1)^2>V_*\), contradicting
\eqref{eq:sf-anchor}. Thus \(j_i<J<W\). Reduction modulo \(A-1\),
using \(Y_j(A)\equiv j\pmod{A-1}\), shows that
\[
 0\le\sum_i(j_i-k_i)^2<W^3<A-1
\]
is divisible by \(A-1\), hence zero. Every index is authenticated.

Conversely let \(m=\max k_i\) and \(z=Y_m(A)\).
If \(J\ge2m\), then
\[
 Y_J(A)\ge2X_m(A)Y_m(A)>2(A-1)z^2.
\]
The sum of actual squared errors is at most \(r z^2\); since
\(2(A-1)-r>J\), the numerator defining \(Q\) is positive and divisible
by \(A-1\). This yields its natural value.

The scale rows are
\[
 L^2-((W+1)^2-1)(WZ)^2=1,\quad Z=z_0+1,\quad A=L+\sigma.
\]
Their positive Pell index is divisible by \(W\), so
\(A\ge X_W(W+1)\ge(W+1)^W\). Nonanchor seeds retain natural
consecutive coordinates \(Y,Z'\) with
\(Y^2+(Z')^2-2AYZ'=1\). Each has an output incidence
\(bY=Z'+R\), \(R\ge0\), \(1\le b<A\). Thus \(Y>0\),
\(Z'<AY\), and \(X=AY-Z'>0\) gives the ordinary principal norm.
The anchor argument therefore applies to these ordinates too.

For every exact request the authenticated pair satisfies
\[
 X_k(A)+(b-A)Y_k(A)=m+(2Ab-b^2-1)h.
\]
The Pell congruence gives \(m\equiv b^k\pmod{2Ab-b^2-1}\).
The frame has \(b,k,m<W\), while its lower scale makes the modulus
larger than \(W^W+W\); hence \(m=b^k\). Canonical positive Pell
residues and their least nonnegative residues give nonnegative quotients.
The common-root equation
\[
 2A(q+1)=(q+1)^2+1+Mh_0,\qquad M=Nq^x,
\]
gives the analogous congruence for all modular numerator requests.
Lemma~\ref{lem:sf-bootstrap} then decodes the first binomial, and only
then Lemma~\ref{lem:sf-fixed} supplies \(F=t!\).

\subsection{Local bounds, soundness, and constructive completeness}
The endpoint equalities now imply
\(R=2+t+x+y+F<2x\) and \(x>2^{32}\).
The retained literal local majorant has \(E=2^{23}\).
At \(x=2^{32}\), \((2x)^E<2^x\); the ratio bound
\[
 (1+1/x)^E<\frac1{1-E/x}<2
\]
propagates the inequality by induction. Thus each required local source
value is \(<R^E<2^x\). The retained label majorant has
\(\ell=2^{22}\), with \(\ell(t+1)<F<x\) for every padded
\(t\ge12\). These are unbounded estimates, not numerical search
cutoffs. The strict beta quotient ranges \(u,v<x\) are complete because
their canonical denominators are at least two.

Every independent representative is \(<S\), so its binomial
coefficients through lower index \(x\) are \(<S^x<q\).
The derived label has \(J\le S^\ell\), making all its needed
coefficients through \(t+1\) also \(<S^x\). All eleven modular requests
are therefore carry-free. For each row \(i-1\), choose \(p\mid d_i\)
and \(m=v_p(d_i)\). The retained one-direction range lemma
\[
 p^e\mid\binom ak\ \Longrightarrow\ a\bmod p^e\in[0,k-1]
\]
recovers the small local values modulo \(p^{xm}\). The index row
\eqref{eq:sf-index} distinguishes \(i-1\). The nonnegative local
source value is \(<2^x\le p^{xm}\), so its zero residue forces its
actual value to vanish. The derived label is similarly bounded.
The unchanged beta and Horn-flow arithmetic yields a finite C222
refutation, and the named retained calibration gives ZFC refutability.

For completeness start with a finite refutation, expand to its proof tree,
and pad \(t\ge12\) beyond every local parameter and label. Use the
inherited beta construction for \(x,y\). Set
\(F=t!\), \(n=F^2-1\), \(B=\prod_i d_i\), \(N=B^x\).
Lemma~\ref{lem:sf-fixed} supplies the ordinary Euclidean remainder
and complementary slack in \eqref{eq:sf-certificate}.
Choose the nine positive independent representatives modulo
\(d_i^{2x}\), initially in \((0,N^2]\), and increase only \(A_a\)
by \(N^2\). All local polynomial congruences to precision \(2x\)
are preserved. A falling-factorial numerator has a term of valuation
at least \(2xm\), while division by \(k!\) loses at most
\(k-1\le x-1\); hence the required binomial is divisible by
\(p^{xm}\). This genuine precision argument supplies range
completeness; no false converse of the range lemma is used.

The nominal indices \(t,y+1,x,n\) are at most
\(x^2+x<2^x\le N\); every independent representative is at most
\(2N^2\). The Cantor-derived label obeys
\[
 J\ge2(A_a+N^2)^2>2N^4>4N^2\ge2\max_i k_i.
\]
Thus the cofinal shift supplies the anchor's completeness hypothesis
without another polynomial row. Choose \(S,q\), denominators and
\(W,L\) in that order. Since \(N\mid q\),
\(\gcd(q+1,M)=1\); the linear common-root congruence has gcd one or
two, and in the latter case its right side is even. Choose \(A\)
arbitrarily large in that solution class. All scale and root quotients
can be natural. Canonical Pell pairs, fresh private anchor witnesses,
\(Q\) from \eqref{eq:sf-anchor}, and Euclidean digit witnesses
complete the graph. The choices are acyclic.

\subsection{A reversible row-basis improvement}
The final balanced endpoint further replaces the 179 precursor residuals
by 42 explicit elementary shears. Each operation, recorded as
\((i,j,k)\) with \(i\ne j\), is
\[
 R_i\longleftarrow R_i+kR_j,\qquad k\in\Z\setminus\{0\}.
\]
Its inverse is the same operation with \(-k\), applied in reverse order.
Thus the old and new residual vectors vanish on exactly the same tuples.
Their sums of squares need not have equal values, but have identical zero
sets over \(\N^{229}\). No coordinate is changed, projected, or made signed.
The resulting quartic has 5,607 collected monomials and coefficient height
65,892,965. The precursor's 6,596 terms and height 83,879,378 are retained
as a control in \texttt{projected229/}.
The complete operations, target residuals and exact metrics are in
\texttt{row\_basis\_certificate.json}; \texttt{check\_row\_basis.py}
checks every forward and inverse row identity, every collected coefficient,
coordinate liveness, and the refreshed embedded endpoint metrics. It rejects
six mutations, including deliberately stale metric metadata.

\subsection{Exact export, natural restoration, and reproducibility}
The independent affine checker validates all 613 defining gates,
392 nonnegative affine restorations, all 285 exported source rows,
and all 4,606 raw coefficients. The balanced precursor contraction has 106
eliminations; both 228-variable variants have 107. Each restoration
is an integer polynomial proved nonnegative by exact rational affine
squares plus nonnegative monomials. The independent checker expands
these certificates, substitutes them back, checks every degree bound,
and independently collects the complete quartic. The literal TeX
residuals are parsed back and compared term by term.

\texttt{verify.py -{}-rebuild} checks all endpoints without mutating them
and rebuilds the selected mathematical files in temporary directories.
The portable source, exact source JSON, retained ordinary proofs and
named Lean calibration inputs are included. Source-path hashes and
modifications are exposed in \texttt{SOURCE\_MAP.json}; the full
unbounded new argument is also available as \texttt{PROOF.md}.
The supplied calibration inputs do not constitute a newly rebuilt full
transitive Lean development or an encoded object-language ZFC proof.

\section{Ten-positive equations by denominator engineering}\label{sec:compact}
This section supplies six complete ordinary integer polynomials, each in ten
strictly positive integer unknowns. Four give useful local operation/degree
tradeoffs; two are retained as dominated controls. All fixed coefficients are
the \emph{whole} indexed BFS943 native tuple, supplied literally in
\texttt{compact/coefficients/}. Its finite reconstruction is
\texttt{rebuild\_coefficients.py}; its largest port has 592,884 bits.
These are not coefficient-free formulas, and fixed ports are not unknowns.

\begin{center}
\begin{tabular}{llrr}
Chart & Weight & Binary operations & Exact degree\\\hline
$Y=q^3s,\ D=XY$ & $V^2$ &77&626\\
$Y=q^3s,\ D=XY$ & $bV$ &78&464\\
$Y=q^3s,\ D=q^4$ & $V^2$ &78&622\\
$Y=q^3s,\ D=q^4$ & $bV$ &79&460\\
$Y=q^2s,\ D=q^3Y$ & $V^2$ &79&518\\
$Y=q^2s,\ D=q^3Y$ & $bV$ &80&384
\end{tabular}
\end{center}
The $78/622$ and $79/518$ points are dominated by $78/464$; they are
uniform-family controls, not frontier claims. The incoming unscaled
$78/630$ and $79/468$ points used $D=Y^2$. The new $77/626$ and $78/464$
points respectively improve both metrics of those matched predecessors.
The $80/384$ point also improves the incoming unshifted $83/436$ point
in both metrics. Older lower-degree constructions remain relevant.
No global optimum or literature-wide priority is claimed.

\paragraph{Attribution and scope.}
The incoming \emph{Certificate Reductions} argument already removed the
artificial scale and recovered a Pell index from a known divisor before
recovering a whole denominator. \emph{Semantic Signs} already proved
unshifted transport and the exclusion of both signed third auxiliary
branches. Pell doubling, product-square extraction, and quotient
normalization are inherited. Here the new choices are $D=XY$, $D=q^4$,
and the partially normalized $Y=q^2s,D=q^3Y$, whose fractional carrier
numerator reuses $Y^2$. The proof below composes explicitly with the
retained ordinary Pell/source theorems. It is not a new end-to-end
proof-assistant certificate and proves neither consistency nor
unsolvability separately.

\subsection{The complete polynomials}
The only independent coordinates are
\[
 J,F,\alpha,i,T,w,s,n,\eta,\zeta\in\Z_{>0}.
\]
The fixed positive integers $\beta,K_0,M_C,M_F$ are the unchanged indexed
tuple, with $\beta\ge31$ odd and $0<M_C<\beta<M_F<2\beta$.
Define polynomial abbreviations
\[
\begin{gathered}
q=\beta J+1,\quad h=q-1,\quad C=q-F-\alpha,\quad X=wq,\\
k=\eta+\zeta,\quad c=kY+\eta,\quad A=Y(X+1)+1,\\
\Delta=A^2-1,\quad H=4A-5,\quad v=XY^2,\\
E=1+v(v+1)k^2,\\
R=[h(q-F)+\alpha](q^2-1)+(M_C+qM_F)J,\\
e=k-R-1,\quad B=(K_0+w)C-F,\\
U=ic^2,\quad a=\Delta U^2E,\quad F_0=1+a,\quad G=EF_0,\\
F_{\rm sq}=1+4aF_0,\quad Q=\Delta(F_{\rm sq}-1),\\
b=cT-R,\quad V=F_{\rm sq}b-c,\quad K=X(3c-2X)+2.
\end{gathered}
\]
Choose $(Y,D)$ from $(q^3s,XY),(q^3s,q^4),(q^2s,q^3Y)$ and choose
$L=V^2$ or $L=bV$. Put $M=hDH$. In the first two charts set
\[
 Z=hD\{H(n-LB)+K\}-H(DB+he).
\]
In the partial chart $Y=q^2s,D=q^3Y$ set instead
\[
 Z=hD\{H(n-LB)+K\}+H\{h(Y^2-e)-DB\}.
\]
In every case the complete polynomial equation is
\begin{equation}\label{eq:newcompact}
 P=Q(Z^2-M^2GV^2)+(M^2G-Z^2)=0.
\end{equation}
Every operation is addition, subtraction or multiplication; no division,
root, exponent variable or hidden quantified intermediate occurs.
For proof only let $z=Z/M$, $\kappa=K/H$, $\rho=e/D$,
$t=B/h$, $W=1+hL$, and let $\sigma=0$ in the cubic charts and
$\sigma=s/q$ in the partial chart. The exact rational identity is
\begin{equation}\label{eq:compactcarrier}
 z=n+\kappa+\sigma-Wt-\rho.
\end{equation}
In the partial chart $D\sigma=Y^2$ is the crucial reuse identity.

\begin{theorem}\label{thm:compact-six}
Each of the six polynomials above has a strictly positive integer zero
if and only if the retained syntactic calculus has a finite ZFC
refutation. Consequently, by the retained whole indexed calibration,
ZFC proves that its positive unsolvability is equivalent to
$\operatorname{Con}(\mathrm{ZFC})$.
\end{theorem}
\begin{proof}
We give the changed unbounded arithmetic argument, retaining the stated
universal Pell, ratio and native-calibration interfaces. Write
$\psi_A(r),\chi_A(r)$ for Pell ordinate and trace.

\emph{Roots before signs.}
Positive inputs imply $q,X,Y\ge32$, $c>k\ge2$, $A>4$, $U\ge c^2$,
$E>1$, $a>c$, $Q>c^2$, and $M>0$. The identities
\[
 F_{\rm sq}=(1+2a)^2,\qquad
 Q=4\Delta^2U^2G,\qquad \gcd(E,F_0)=1
\]
are unconditional. For any integral $b$, $V\ne0$ and $QV^2-1>0$.
Equation~\eqref{eq:newcompact} gives
$(Q-1)z^2=G(QV^2-1)$. Since $G\mid Q$, both other integer factors
are $-1$ modulo $G$; prime valuations force $G=g^2$.
Coprimality gives $E=\tau^2,F_0=f_0^2,g=\tau f_0$ with positive roots.
The rational number $y=|z|/g$ satisfies
\begin{equation}\label{eq:adjacent-new}
 QV^2-(Q-1)y^2=1.
\end{equation}
Rational adjacent-Pell descent makes $y$ integral: if its reduced
denominator is $m$, then $m^2\mid Q-1$; the norm-preserving map
\[
 (v,y)\longmapsto((2Q-1)v-2(Q-1)y, (2Q-1)y-2Qv)
\]
keeps the first coordinate integral, keeps the reduced denominator,
and strictly decreases $|v|>1$ after sign reflection. At $|v|=1$,
$y=1$. Thus $z=\varepsilon gy$ is integral.

Put $f=1+2a,I=2Ug$. Then
$f^2=F_{\rm sq}$, $f^2-\Delta I^2=1$, $c^2\mid I$ and
$Q=(\Delta I)^2$. The retained signed all-index folding theorem,
which allows either sign of $V$, applies to these equations and
$V=-c\pmod f$. It gives $c=\psi_A(p)$ for an odd $p\ge1$.
The first norm gives $k=2\psi_{2v+1}(n_0)$, $n_0\ge1$.
The value $n_0=1$ would give $c=2Y+1$ strictly between
$\psi_A(1)=1$ and $\psi_A(2)=2A$. Hence
\begin{equation}\label{eq:earlyc}
 k\ge8XY^2+4,\qquad c>kY>8XY^3\ge8q^4,
 \qquad \kappa>0.
\end{equation}
No sign of $B,C,R,b$ has been used.

\emph{Sign-free transport.}
At a zero, $b=0$ would make
$y^2=c^2+(c^2-1)/(Q-1)$ lie between consecutive integers. Thus
$\operatorname{sgn}V=\operatorname{sgn}b$ and $|V|>|b|$.
Also $a^2-G=E[\Delta U^2E(\Delta U^2-1)-1]>0$,
$F_{\rm sq}-c>3a>3g$, and $y<3|V|/2$.
Therefore $gy<V^2/2$.

For $L=bV$ a stronger bound is needed. Suppose $p=3$, so
$c=4A^2-1$. If $n_0\ge3$, then $kY>32X^2Y^5$, whereas
$c<16X^2Y^2$, a contradiction. Hence $n_0=2$ and $k=8XY^2+4$.
For the integer $u=X+2-2Y$ one obtains
\[
 \eta=4Y^2(Xu+1)+8Y(X+1)+3-4Y.
\]
If $u\le-1$ this is negative, since
$Y(X-1)>2(X+1)$ for $X,Y\ge32$.
If $u\ge2$ it exceeds $k$. Hence $u\in\{0,1\}$, or
$2Y-X\in\{1,2\}$, impossible since $q\mid X,Y$.
Thus $p\ge5$ and $c>c_3=4\Delta+3$.

Writing $S=\Delta I$, Pell classification of
\eqref{eq:adjacent-new} gives $|V|=\chi_S(j)/S$ for odd $j$.
The first possibility is $|V|=1$, already excluded. For the third,
$V=\pm(4Q-3)$ gives respectively
\[
 F_{\rm sq}(b-4\Delta)=c-c_3,\qquad
 F_{\rm sq}(b+4\Delta)=c+c_3.
\]
Both right sides lie strictly between $0$ and $F_{\rm sq}$.
Thus $j\ge5$, $|V|>8Q^2$, and, since $Q>F_{\rm sq}>2c$,
\[
 |b|>\frac{8Q^2-c}{F_{\rm sq}}>4F_{\rm sq}>4g.
\]
Consequently $L=bV=|b||V|>0$ and $gy<3L/8<L/2$.

For either weight $e<b$ because $b-e=cT-k+1>0$, so
$\rho<|b|/D$. If $B\le-1$, then $Wt<-L$; the carrier gives
\[
 n< L/2-L+|b|/D<0,
\]
using $D>2$, $|V|>|b|$ and $|V|>1$. Contradiction.
Thus $B\ge0$, and $F>0$ forces $C>0$.
Only now use the mask estimate $3q<R<q^4-2$.
Equation~\eqref{eq:earlyc} gives $c>2R$, hence $b>0,V>0$.
The retained normalized-rank and signed-index theorems, using
$c^2\mid I$, $c>2R,c>2p$ and the fixed-minus residues, give $p=R$.
Their hypotheses do not require the removed shift $T+\alpha$.
Odd $R$ then gives even $q$ and odd $J$.

\emph{Recover a small index before the full denominator.}
The recurrence at $P_0=2XY^2+1$ gives
\[
 k\equiv2n_0\pmod {4XY^2}.
\]
Since $P_0>A$, $n_0\ge R$ would imply $k/2\ge c$.
Thus $1\le n_0<R$. Integral $z$ gives $D\mid hHe$ in the cubic
charts, where $q^4\mid D$. Since $H=-1\pmod q$ and $h=q-1$,
this yields $q^4\mid e$.
In the partial chart clearing $qhH$ gives $D\mid qhHe$;
$q^5\mid D$ again gives $q^4\mid e$.
Crucially $q^4\mid4XY^2$ in \emph{all} these charts.
Therefore $q^4\mid2n_0-R-1$, and
$|2n_0-R-1|<R<q^4$ forces $2n_0=R+1$.
The full recurrence now gives $4XY^2\mid e$.
Each selected $D$ divides $4XY^2$; for $D=q^3Y$ the quotient is
$4ws$. Thus $\rho$ is integral. No false coprimality of $Y$ and $h$
was invoked.

\emph{Projection and source decoding.}
Write $R=2r+1$, so $r\ge48$. The retained ratio theorem at scale one,
with $\xi=(X+1)^{2r}/X^r$, gives
\[
 \xi<\frac{\psi_A(2r+1)}{\psi_{P_0}(r+1)}
 <\xi\left(1+\frac1{Y(X+1)}\right)^{2r}.
\]
Together with $Y<c/k<Y+1$, this gives $Y>X^r/3$ and
$X^{r+1}<H$. With $\rho$ integral the carrier gives
$H\mid hK$ in the cubic charts, $H\mid qhK$ in the partial chart.
Either implies the inherited premise $H\mid q^3hK$.
Universal projection and odd $H$ give
\[
 H\mid q^3h(X-2^R)(2^RX+1).
\]
The retained whole-product bound puts the absolute value of this
\emph{entire} product below $X^{r+1}<H$, hence $X=2^R$.
In particular $\kappa$ is integral; no individual nonunit factor was
cancelled modulo $H$.
The cubic carrier gives $h\mid B$. The partial carrier gives
$s/q-WB/h\in\Z$; coprimality of $q,h$ gives $q\mid s$ and $h\mid B$.
Thus $q^3\mid Y$ and $t\ge0$ is integral in every chart.
The retained fixed-minus parity and half-binomial rounding theorems
now have their full hypotheses. Their fresh reconstruction at
$A_{\rm source}=Y(X+1)+2$, followed by the whole indexed native
calibration, supplies a finite ZFC refutation.

\emph{Positive converse.}
A refutation supplies the retained canonical outer data,
$X=2^R$, $R\equiv3\pmod4$, $R>3q$, $q^3\mid Y$, integral $t\ge0$,
and the required half-binomial $Y$. Set $s=Y/q^3$ in cubic charts,
$s=Y/q^2$ in the partial chart. Form the scale-one bases and canonical
$c=\psi_A(R)$, $k=2\psi_{P_0}((R+1)/2)$, $\tau=\chi_{P_0}((R+1)/2)$.
The ratio/rounding theorem gives $0<\eta=c-kY<k$, so
$\zeta=k-\eta>0$. Also $4XY^2\mid e$, hence $\rho\ge0$ is integral,
and projection gives integral $0<\kappa<Xc$.
Choose a positive matrix-period multiple at $A$ modulo $c^2\tau$,
with seed pair $(f_0,I_0)$ and $c^2\tau\mid I_0$.
Put $i=I_0/(c^2\tau)>0$ and double the pair; this produces exactly
$f=2f_0^2-1$, $I=2f_0I_0$, and all literal root definitions above.
The retained auxiliary completeness theorem supplies positive $V,y$ and
\[
 T=\frac{V+c+Rf^2}{cf^2}\in\Z_{>0}.
\]
No subtraction of $\alpha$ is needed. Set
$n=gy+Wt+\rho-\kappa-\sigma$.
In the partial chart $\sigma=Y/q^3$ is now integral; in cubic charts
it is zero. Since $\sigma<Y<Xc$, $\kappa<Xc$, and
$f_0>(A-1)I_0>2Xc$, we have $g>\kappa+\sigma$ and therefore $n>0$.
The carrier is $z=gy$, proving~\eqref{eq:newcompact}. All independent
coordinates are positive; this is a fresh witness construction, not an
unjustified positive inverse shear.
\end{proof}

\subsection{Exact degree and reproducibility}
Let $q_0=\beta J$, $k_0=\eta+\zeta$, $C_0=q_0-F-\alpha$.
The full highest homogeneous forms are respectively
\begin{align*}
 (XY,V^2):&\quad2^{14}q_0^{378}w^{48}s^{110}i^{20}k_0^{64}T^4C_0^2,\\
 (XY,bV):&\quad-2^{10}q_0^{284}w^{36}s^{82}i^{14}k_0^{46}T^2,\\
 (q^4,V^2):&\quad2^{14}q_0^{378}w^{46}s^{108}i^{20}k_0^{64}T^4C_0^2,\\
 (q^4,bV):&\quad-2^{10}q_0^{284}w^{34}s^{80}i^{14}k_0^{46}T^2,\\
 (q^3Y,V^2):&\quad2^{14}q_0^{272}w^{46}s^{110}i^{20}k_0^{64}T^4C_0^2,\\
 (q^3Y,bV):&\quad-2^{10}q_0^{206}w^{34}s^{82}i^{14}k_0^{46}T^2.
\end{align*}
For $V^2$ the unique top summand is $QZ^2$; for $bV$ it is
$-QM^2GV^2$. The displayed nonzero forms prove the claimed exact global
degrees at the actual positive $\beta$. Every binary operation is charged,
including squares and multiplication by scalar literals. All inputs,
ports and gates are live. Six deterministic DAGs, frozen parent data,
coefficient bytes, recipe and mathematical provenance accompany the proof.
The independent checker covers every row by exact polynomial-cut identities,
propagates complete leading forms, authenticates all fixed bytes, and rejects
48 deliberate algebraic corruptions. Normal and optimized Python replays
and deterministic rebuilds are supplied. These finite checks are not
substitutes for the unbounded theorem above.

A tempting $Y=qs,D=q^6$ route at $81/410$ or $82/304$ was rejected:
$4XY^2$ is only visibly divisible by $q^3$, insufficient for its index
argument. It is not an accepted equation or delivered endpoint.

\section{Complete fixed-coefficient recipe}\label{sec:compact-constants}
The following finite recipe is retained from the inspected ProveIt indexed
constant specification, whose exact source and SHA-256 are recorded in
\texttt{compact/provenance/indexed\_constants.tex} and
\texttt{compact/SOURCE\_PINS.json}. It defines the entire unchanged
BFS943 tuple, not a new specialization. The symbol $K$ in the inherited
recipe is renamed $K_0$ to agree with Section~\ref{sec:compact}; symbols
in this section are local fixed data, not additional unknowns. The supplied
standard-library reconstruction checks every resulting byte against all
six literal circuits. The source shift is already included in $M_F$.
\begingroup
\providecommand{\fl}[1]{\left\lfloor#1\right\rfloor}
\providecommand{\rg}[2]{\{#1,\ldots,#2\}}
\textbf{The four fixed integers.}
Every index, set, sum, product and floor below defines a constant, not
another unknown. Put
\[
\rho=2^{67},\quad B=\rho^{8807},\quad\beta=B-1,\quad
G_m(x)=\sum_{j=0}^{m-1}x^j.
\]
Define
\[
\begin{split}
\mathcal E={}&\rg0{942}\cup\rg{983}{992}\cup\rg{1003}{1011}\\
&\cup\rg{1013}{1031}\cup\{1063,1064,1095,2195,3296\}.
\end{split}
\]
Then
\begin{align*}
M_C={}&B-1-\sum_{e\in\mathcal E\setminus\{1\}}\rho^e-2\rho^{1063},\\
M_F={}&B-1+4+(2^{51}-2)(\rho^{4393}+\rho^{5495})\\
&+\rho^{1013}G_{10}(\rho)+\rho^{1065}G_9(\rho)
  +\rho^{2195}+\rho^{3296}\\
&+(2^{65}-1)\rho^{8792}G_{13}(\rho)+(2^{19}-1)\rho^{8805},\\
K_0={}&\rho^{30}+\rho^{62}+\rho^{1100}+\rho^{2201}+\rho^{8849}\\
&+(\rho^{4393}+\rho^{5495})\Phi(1/\rho),
\end{align*}
where
\[
\begin{split}
\Phi(x)={}&(1+2^{51})G_{943}(x)
 -2^{10}\sum_{i=0}^{942}v(t_i)x^i\\
&+\sum_{(b,e,m)\in\mathcal C}2^b x^e G_m(2x),\\
\mathcal C={}&\{(10,1023,8),(18,1013,9),(27,1003,8),\\
&(35,983,9),(44,1031,1),(45,1022,1),(46,1011,1),\\
&(47,992,1),(48,1095,1),(49,2195,1),(50,3296,1)\}.
\end{split}
\]
The rational intermediate in $K_0$ cancels exactly; all four results are positive integers.

\textbf{The 943 indexed fixed tiles.}
Use the actions $u_j,v_j,d_j$ decoded below. Let
\[
\begin{gathered}
h(q,b)=2(q+1)+b,\quad
\mathcal R=\{u_j:d_j=1,\ 0\le j<430\},\\
\mathcal L=\{u_j:d_j=0,\ 0\le j<430\},\\
r_q=6+|\mathcal R\cap\rg0{q-1}|,\quad
l_q=205+|\mathcal L\cap\rg0{q-1}|.
\end{gathered}
\]
An interval ending at $-1$ is empty. The table defines $t_i$ for
$0\le i\le942$; all index ranges are fixed:
\[
\setlength{\arraycolsep}{1.4pt}\begin{array}{c|l}
\text{index }i&\text{tile }t_i=(w,n,e,s)\\\hline
0,1&(2,435,3,1),\ (5,0,0,0)\\
j+2\ (j=2q+b<430)&((1-d_j)l_{u_j},h(q,b),d_jr_{u_j},v_j)\\
432+b\ (b=0,1)&(0,h(215,b),0,h(215,b))\\
434+b\ (b=0,1)&(0,b,0,b)\\
424+2r_q+b\ (q\in\mathcal R,b=0,1)&(r_q,b,0,h(q,b))\\
424+2l_q+b\ (q\in\mathcal L,b=0,1)&(0,b,l_q,h(q,b))\\
924+j\ (0\le j<4)&(0,j+430\fl{j/2},0,438)\\
928&(2,434,2,0)\\
929&(3,436,4,4)\\
930&(4,436,4,0)\\
497+t\ (t=434,435,436)&(1,438,1,t)\\
\substack{934+3a+b-6\fl{a/4}\\a\in\{0,1,4\},\ b=0,1,2}&(a,437,b,437)
\end{array}
\]
For $t=(w,n,e,s)$ put
\[
v(t)=w+2^8n+2^{17}e+2^{25}s+127\,2^{34}.
\]

\textbf{A closed decoder for the 430 actions.}
Put $D_b(A,k)=\fl{A/b^k}\bmod b$, where
$a\bmod b=a-b\fl{a/b}$. For $0\le q<215$, $b=0,1$, $j=2q+b$:
\[
\begin{gathered}
N=\mathcal A\bmod2^{636},\quad
T_0=\fl{\mathcal A/2^{636}}\bmod2^{1675},\\
E_0=\fl{\mathcal A/2^{2311}},\\
\nu_j=D_2(N,j),\quad \sigma_j=\sum_{k=0}^{j-1}\nu_k,\\
m_q=D_2(N,430+q),\quad \tau_q=\sum_{k=0}^{q-1}m_k,\\
F_q=14(1-m_q)+m_qD_{16}(E_0,\tau_q),\\
u_j=\nu_j(1+\sigma_j)+(1-\nu_j)D_{216}(T_0,j-\sigma_j),\\
v_j=D_2(F_q,2b),\quad d_j=D_2(F_q,2b+1).
\end{gathered}
\]
Empty sums are zero. Here $(v_0,d_0,u_0)=(1,1,1)$; the constructor tile
uses this actual first target, including its state renaming.

\textbf{The integer $\mathcal A$.}
The alphabet, in order of digit values $0,\ldots,61$, is\par
{\fontsize{9}{10}\selectfont\ttfamily 0123456789ABCDEFGHIJKLMNOPQRSTUVWXYZabcdefghijklmnopqrstuvwxyz}.\par
Concatenate the following 436 characters as digits $e_0,\ldots,e_{435}$;
put $\mathcal A=\sum_{j=0}^{435}e_j62^{435-j}$.\par
\begingroup\fontsize{9}{10}\selectfont\ttfamily
BxEnhvy6NipzkWmROo0GsA3uK3twWLxscLq20UTumy5Cdu461LwD6g7RDNzF8Y\par
iYDFTpB8ehAypd7d6OsOL4APamnJeNXQypUezCmW9n4oXP6XfgmiWuAWZ8ZJ7U\par
OP7gu4kVdlGZqZHULQMar1dXNeJW4i7oNVI8usp8BlovjjZbwPCovZI328fRUq\par
4q92R81s5GjWQaQfGhkq9UYMK3rRMaBmsOIBkJqc2cj2VRIZ6mjzkBIOkW9SjZ\par
FzeA6J5u4j3sBL3SdgwXzM9mHvQivu8EJ1TBTpRhHDqZIbmYgwQ8VqTAFotl6F\par
We1YC554GKrVHsA8oUGKSA6F6k5KREBSD0Tgd7bFALY84LLbpw574BQyvdoHrl\par
0Oe35uwDf15DpT6VJWZBjSbHPpPAINYTte0eUxnTFTKBlGj8Am6sb1gLldh8WZ\par
8P\par
\endgroup

\endgroup

\section{Primitive odd carriers: fourteen integers at degree 86}
\label{sec:root-primitives}

The auxiliary Pell root and the triangular positivity coordinate can be
combined directly, without carrying the first Pell root through a three-root
norm. The resulting complete equation has fourteen unrestricted integer
unknowns, exact degree $86$, and $112$ live binary operations. It improves
both degree and arithmetic length over the retained $14/88/119$ point.
A $14/94/110$ companion saves two more operations at greater degree.
The proof changes no fixed numeral or native/compiler calibration.

\begin{center}\small
\begin{tabular}{lrrr}
Construction & Integer unknowns & Exact degree & Operations\\\hline
Primitive auxiliary/guard carrier &14&86&112\\
Primitive strong/guard carrier &14&94&110\\
Exact recompilation of old two-root equation &14&88&115\\
Exact recompilation of old three-root equation &13&156&125
\end{tabular}
\end{center}
The $88/115$ row is dominated by $86/112$ and is retained to isolate the
compiler-only saving. Arithmetic operations mean live binary addition,
subtraction and multiplication, including multiplication by constants;
these counts do not measure bit complexity.

\subsection{Common arithmetic front and the precise inherited interface}
The six fixed integers are exactly the retained source34 tuple
$(\beta,K_0,M_C,M_F,A_c,B_c)$, with $M_F$ already shifted. In particular
$\beta=2^{1064921}-1$, $A_c=2^{1113}$ and $B_c=2^{666104}$.
Their complete authenticated payloads and finite reconstruction recipe
accompany the DAGs. Define polynomial abbreviations
\[
\begin{gathered}
J=6j+1,\quad q_m=\beta J,\quad q=q_m+1,\\
W=12w+4,\quad X=qW,\quad Y=v,\\
F=(K_0+W)C-zq_m,\\
R=(q(q-F)-C)(q^2-1)+(M_C+qM_F)J,\\
E=XY,\quad a=E+Y,\quad H=4a+3,\quad\Delta=a^2+H,\\
k=R+1+hE,\quad\eta=c-Yk,\quad L=E(Yk),\\
I=(2r+1)c^2,\quad M=2q^3,\quad U=MH,\\
K_\ell=X(2X-3c)-2,\quad Z=Un+MK_\ell+H(Y-q^3),\\
D_e=(A_cF-q)^2+(B_cC-q)^2,\\
\mathcal B=(Yk)^2(q^2-D_e)-q^2\bigl(4(Y\eta)^2+(kq^3)^2\bigr),\\
A=(Z+Uk)^2+4U(LZ-U).
\end{gathered}
\]
All names in this display are polynomial abbreviations, not additional witnesses.
The fifteen-integer parent additionally uses $f,y,p_1,p_2,p_3$ and
\begin{align*}
 Q&=\Delta(f^2-1),&V&=f^2(cT-R)-c,\\
 B_s&=f^2-\Delta I^2-1,&
 C_a&=V^2-(Q-1)y(2V+y)-1,\\
 B_g&=\mathcal B-1-p_1^2-p_2^2-p_3(p_3+1).
\end{align*}
Its equation is $A^2+B_s^2+C_a^2+B_g^2=0$.
The retained ordinary proof identifies integer-zero existence with finite
syntactic ZFC refutability. The following new maps restore an actual zero of
this precise parent, rather than introducing a new source interpretation.

The facts needed before extraction hold on every integer assignment.
The integer $q$ is even and nonzero, $H$ is odd and nonzero, and $U\ne0$.
Also $|X|\ge8$, so $a=(X+1)Y$ is zero or has $|a|\ge7$; hence
$\Delta=(a+2)^2-1\ge3$. The polynomial $A$ is a square modulo $4$,
so $A\ne-2$. Since $\Delta$ is $0$ or $3$ modulo $4$, $B_s$ is $0,1$
or $3$ modulo $4$, and $B_s\ne-2$ as well.

\subsection{An odd carrier with no exceptional zero branch}
\begin{lemma}[Primitive merger]
Let $b,d,e,s\in\Z$, with $d\ne0$. Set
\[
 G=4b+1,\qquad K=d\bigl(s(s+1)-b\bigr)-e,
 \qquad N=K^2-Gde.
\]
Then $N=0$ if and only if there are integers $g,u$ with
\[
 g^2=G,\qquad du^2=e,\qquad 2s+1=g+2u.
\]
\end{lemma}
\begin{proof}
Put $x=2s+1\ne0$. Direct expansion gives
\[
 16N=\bigl[d(x^2+G)-4e\bigr]^2-4d^2x^2G.
\]
At a zero, $g=[d(x^2+G)-4e]/(2dx)$ is rational with integral square $G$,
so it is an integer. Since $G\equiv1\pmod4$, $g$ is odd. Thus
$u=(x-g)/2$ is integral, and the defining relation for $g$ gives $du^2=e$.
Conversely substitution gives $K=dgu$, hence $N=0$.
The argument permits negative $d,e$ and proves, rather than assumes,
nonnegative $G$. In particular it loses no irrational-cancellation or
zero-carrier branch.
\end{proof}
The generic primitive identity is already present in the retained parity-norm
literature; the new result is its placement in this free-scale interface and
the complete degree/cost reduction.

\subsection{The complete degree-86 equation}
Use the fourteen unrestricted inputs
\[
 (j,C,z,h,r,v,w,T,c,n,f,p_2,p_1,s).
\]
Keep the front and $Q,V$ above, and define
\begin{align*}
 B&=\mathcal B-p_1^2-p_2^2,&G&=4B-3,&\delta&=Q-1,&C_y&=QV^2-1,\\
 K&=\delta\bigl(s(s+1)-B+1\bigr)-C_y,&
 N&=K^2-G\delta C_y,&S&=f^2-\Delta I^2.
\end{align*}
The single equation is
\begin{equation}\label{eq:root86}
 \boxed{\;P_{86}=(A+1)S(N+1)-1=0.\;}
\end{equation}
This is an ordinary polynomial with fixed natural powers and no divisions,
radicals or hidden witnesses. Its full division-free DAG is supplied.

\begin{proposition}
Equation~\eqref{eq:root86} has an integer zero if and only if the fixed
fifteen-variable parent does.
\end{proposition}
\begin{proof}
At a zero, the three integer factors $(A+1),S,(N+1)$ are units. The first
cannot be $-1$ because $A\ne-2$, and the second cannot be $-1$ because
$S=B_s+1$ and $B_s\ne-2$. Thus $A=B_s=N=0$.
The strong row implies $Q=\Delta^2I^2$, which is $0$ or at least $9$;
therefore $\delta\ne0$. Apply the lemma with
$(b,d,e)=(B-1,\delta,C_y)$, obtaining odd $g$ and integral $u$.
Restore
\[
 p_{3,\mathrm{old}}=(g-1)/2,\qquad y_{\mathrm{old}}=u-V.
\]
All other inputs are unchanged. Then $B_g=0$ and
\[
 C_{a,\mathrm{old}}=QV^2-(Q-1)(V+y_{\mathrm{old}})^2-1
                  =C_y-\delta u^2=0.
\]
This includes $I=0$, where $Q=0$, $\delta=-1$ and $C_y=-1$.

Conversely, from any parent zero set $s=p_{3,\mathrm{old}}+V+y_{\mathrm{old}}$
and retain all other common inputs. With
$g=2p_{3,\mathrm{old}}+1$ and $u=V+y_{\mathrm{old}}$, the lemma applies;
hence $N=0$ and the other two factors are $1$.
These maps work on every zero, not just a canonical family.
\end{proof}

\subsection{The degree-94, 110-operation alternative}
Use inputs $(j,C,z,h,r,v,w,T,c,n,y,p_1,p_2,s)$ and replace $f^2$ by
\[
 D_s=\Delta I^2,\quad F_s=1+D_s,\quad Q=\Delta D_s,
 \quad V=F_s(cT-R)-c.
\]
Keep $B=\mathcal B-p_1^2-p_2^2$, $G=4B-3$, and put
\[
 K=s(s+1)-B-D_s,\quad N=K^2-GF_s,
 \quad E_a=V^2-(Q-1)y(2V+y).
\]
The complete equation is
\[
 \boxed{\;P_{94}=(A+1)(N+1)E_a-1=0.\;}
\]
Here $F_s$ is $0$ or $1$ modulo $4$, so $N=K^2-GF_s$ is $0,1$ or $3$
modulo $4$ and cannot be $-2$. The first two factors must consequently
both be $1$, forcing $A=N=0$ and $E_a=1$. Apply the same lemma with
$(b,d,e)=(B-1,1,F_s)$ to recover $f^2=F_s$ and odd $g^2=G$; restore
$p_3=(g-1)/2$. Conversely set $s=p_3+f$ at a parent zero. This proves
both directions without additional Pell existence assumptions.

\subsection{Exact degrees, compiler savings, and verification}
Write
\begin{gather*}
 q_0=6\beta j,\quad X_0=12q_0w,\quad Y_0=v,\quad
 R_0=q_0^4z-12q_0^3wC,\quad a_0=X_0Y_0,\\
 U_0=8q_0^3X_0Y_0,\quad L_0=X_0Y_0^2R_0,
 \quad Z_0=U_0n+4q_0^3X_0^2.
\end{gather*}
The complete degree-22 carrier form is
$A_{22}=U_0^2R_0^2+4U_0L_0Z_0$; the degree-18 budget form is
$-R_0^2q_0^8$. The first carrier summand is the already-disclosed correction
to the inherited printed leading form, not a new defect discovery.
Its omission does not change the old degree or the following $n$-dependent
coefficient witnesses. The checker propagates the entire actual DAG.

For $P_{86}$ the norm has degree $52$, with top form
$a_0^4f^4R_0^4q_0^{16}$; the strong unit has degree $12$.
For $P_{94}$ the norm and auxiliary unit both have degree $36$, with top
forms $R_0^4q_0^{16}$ and $32a_0^6r^4c^8R_0y$.
The unique highest product in each finalizer gives
\begin{align*}
 [j^{51}w^9v^{10}z^5f^4r^2c^4n]P_{86}
   &=-2^{79}3^{60}\beta^{51}\ne0,\\
 [j^{55}w^9v^{10}z^6r^4c^8yn]P_{94}
   &=2^{86}3^{64}\beta^{55}\ne0.
\end{align*}
Their exponent sums are $86$ and $94$, respectively, so the degrees are
attained at the actual fixed numeral tuple.

The old $13/156/127$ evaluator computed an auxiliary residual $E-1$ and
later added $1$ back. Compiling $E$ directly removes two operations, giving
$13/156/125$ with the identical polynomial. The old $14/88/119$ evaluator
did the same for both auxiliary and strong residuals; removing both pairs
gives $14/88/115$, again pointwise-identical.

The independent checker binds every named polynomial cut, all live gates,
complete leading forms and actual nonzero coefficients. It checks the
primitive and lifting identities, signed and residue controls, complete
source/numeral bytes and all $453$ separate live-gate corruptions, normally
and under Python optimization. A rebuild reproduces every DAG byte for byte.
The numeral recipe separately reconstructs all six fixed integers.
These are finite algebraic checks, not a decision of zero existence. The
unbounded signed-Pell/native/compiler/refutation calibration remains the
explicit inherited ordinary dependency; no new end-to-end proof-assistant or
object-language consistency certificate is asserted.

\section{Reproduction and logical scope}
The release includes literal equations, complete fixed coefficient data
where used, source provenance, forward and inverse transformations,
reconstruction programs, independent exact checks, and the editable
sources of this report. The top-level replay command is described in
\texttt{README.md}. Normal and optimized Python runs check the same
mathematical obligations without depending on assertions. Full collected
coefficient checks, all-register circuit checks and deliberately corrupted
controls supplement the ordinary proofs.

A finite replay cannot certify an unbounded Pell theorem by sampling.
The mathematical claims instead compose the displayed universal arguments
with the named retained source theorems. The direct-calculus branches use
the numerical C222 refutation interface and its ZF--ZFC consistency
calibration, not a K215 simulation or Jones universal polynomial. The
compact and root-certificate branches retain their distinct native
coefficient/refutation calibrations. No tuple from one branch may be
substituted into another merely because its port names look similar.

No upstream repository was changed. No literature-wide priority,
minimum variable count, lower bound on circuit complexity, or fresh
closed-kernel proof is claimed. Source-reviewed improvements arriving
after the comparison pin do not retroactively change the identities of
the sealed equations or the measured comparisons.
\appendix
\section{An alternate coherent radix55 specialization}
The separate \texttt{compact\_radix55/} component retains the six free-port
DAGs, domains, counts and degrees, but binds \emph{all four} ports to the
accepted current radix55 native instance. The original radix67 component
is unchanged. The specialized polynomials differ: their positive $\beta$
values, and hence their displayed leading coefficients, differ.

This is credited reuse of ProveIt's accepted wrapping-controller encoding,
not a new encoding claim. Total raw coefficient storage decreases from
295,388 to 242,484 bytes (52,904 fewer); the largest port decreases from
592,884 to 486,696 bits. No printed-page or PDF-byte gain is claimed.
The full independent reconstruction and whole binding are supplied.
With the same finite $\Phi$, support and ordered 943 tiles as the fixed
recipe, now let $\rho=2^{55}$, $\mathcal B=\rho^{8807}$ and
$\beta=\mathcal B-1$. Define $K_0,M_C$ by the same formal exponent
recipes, but replace the \emph{entire} source-mask padding by
\[
\begin{aligned}
M_F={}&\beta+4+(2^{51}-2)(\rho^{4393}+\rho^{5495})+C_{\rm core}\\
 &+(\rho-1)\sum_{j=8792}^{8806}\rho^j+(2^{41}-4)\rho^{7691}.
\end{aligned}
\]
where $C_{\rm core}=\sum_{j=1013}^{1022}\rho^j+
\sum_{j=1065}^{1073}\rho^j+\rho^{2195}+\rho^{3296}$.
Simply changing the radix in the old padding would be wrong.

The bridge uses the complete accepted native theorem, not numerical boxes
alone. Its repaired wrapping controller, with the two actual unit residues
3422 and 500242 at 8807 and 581197, constructs fresh canonical words and their
\emph{actual} rank with $R\equiv3\pmod4$, $\operatorname{pc}(R)=3dN+2$,
$q=\mathcal B^N$, $X=2^R$, and $Y/q^3$ positive odd.
Its positive transport coordinate satisfies $B/(q-1)=z_{\rm source}-1\ge0$.
These are exactly the outer data required by the six-chart converse;
all seed, strong and auxiliary Pell witnesses are rebuilt afresh.
For soundness the six-chart arithmetic restores the same half-binomial,
$q^3\mid Y$, rank and transport fractions. Reconstruct the positive
unscaled source at $A_{\rm source}=Y(X+1)+2$, choose its raw positive
auxiliary quotient $T^*$. With the fresh source discriminant
$\Delta_*=A_{\rm source}^2-1$ and fresh source witnesses $i_*,c_*$,
set $S=\Delta_* i_*c_*^2$. These are not the reduction chart's discriminant or private witnesses.
Its accepted P66 factors satisfy
$N_0=N_1=N_a=N_k=N_t=1$, $N_s=\Delta_*$, so that whole radix55
polynomial vanishes and its accepted calibration supplies a refutation.
No old radix67 native word, isolated coefficient or private witness is reused.
The full bridge proof and exact source pins are supplied; the old false
modulo 5 unit is not imported into the repaired controller.

\end{document}
